% ===== §1 =====
\section{Understanding as a Crossing of Worlds}
\label{sec:intro}

\noindent
There is a way of speaking about understanding another person that treats it as the accurate reception of a message. On this way of speaking, the other has thoughts and feelings, encodes them in words, and sends them across, and to understand is to decode what was sent and to recover the meaning intact. The measure of a good bond, on this picture, is the transparency of the channel, the degree to which what one person means arrives in the other without loss. Misunderstanding is noise on the line, and the ideal at the horizon is a communication so complete that two people would be present to one another without remainder.

This paper proceeds from a different account of what understanding is. A companion work developed a generative relational epistemology, on which no single symbolic system carries a concept without remainder, and a concept becomes observable only through the relations among the several symbolic systems that render it. Knowing, on that account, is a generative observation and not the reception of a message, produced in the movement between renderings and enriched across history. The companion work observed, near its close, that a person is the bearer of a symbolic system, a rendering of the world shaped by an experience, a history, a language, a form of life that no other person exactly shares. From this it follows that the meeting of two persons is a crossing of symbolic worlds, of the same kind as the crossing of disciplines, and that understanding another is the generative observation of a matter that neither world holds alone, and not the decoding of a message.

The present paper takes this account into the domain of intimacy. Intimacy is the place where the crossing of two symbolic worlds is at once deepest and most consequential, where the worlds are most entangled and the cost of failing to cross them is most severe. It is also the place where the received picture of transparent communication does its greatest damage, since it sets as the goal of the bond a fusion that the relational account holds to be neither possible nor desirable. What a good intimate bond asks is that each learn to enter the other and to generate, in the relation between them, what neither held before, and not that two worlds become one, an understanding and a value that neither held before.

The paper has a plain structure. Part One recalls the epistemology on which it stands and draws the bridge from concepts to persons. Part Two builds a functional case, setting the work that relational understanding does in an intimate bond beside the established theories and findings that treat trust, shared value, conflict, and the quality of a shared life, and locating the contribution in two characters of the account, the relational and the generative. Part Three then argues that a functional case is not sufficient, since an efficient possession of another can counterfeit the functions of understanding, and it grounds a normative distinction in the ethics of non-possession. Part Four states the problem that sets intimacy apart from other domains of understanding, that its capacities are wagered and not rehearsed. Part Five proposes a model, a set of capacities and their cultivation. Part Six closes without a synthesis, in keeping with the account of understanding the paper develops.

A word on stance belongs here, since it governs what follows. This paper does not claim that a given understanding of one's partner is correct, and the reader who looks for such a claim will not find it. The epistemology on which the paper stands holds that truth is generated in a historical and contradictory movement and is never possessed as a finished correctness, and understanding between persons is subject to the same condition. What the paper asks of an understanding is that it be legitimate, ethically, and not that it be correct, and it holds even this legitimacy to be a historically generated measure and no final guarantee. The distinction between the correctness the paper declines to claim and the legitimacy it does defend runs through the whole of what follows, and it is drawn out in Part Three.
